\documentclass[11pt]{article}
\usepackage[margin=1in]{geometry}
\usepackage{amsmath,amssymb,amsthm,hyperref,enumitem}
\hypersetup{colorlinks=true,linkcolor=blue,urlcolor=blue}
\newtheorem{theorem}{Theorem}
\newcommand{\bb}{\mathcal B}
\title{Positive Bernstein coefficients of Bernoulli densities}
\author{Principia Math In Progress}
\date{21 September 2026}
\begin{document}
\maketitle
\noindent\textbf{MathDB 356321. Full intended theorem; source review pending.}
The identified journal final remains unread, so this package earns
\textbf{zero new campaign credit}. The endpoint wording correction alone
is not counted as a research solution.

\section*{Task 1: the all-degree positivity theorem}
Let \(b_n(x)=B_n(x)/n!\), with Bernoulli numbers \(B_n=B_n(0)\) in the
convention \(B_1=-1/2\). Let
\(\bb_{r,n}(x)=\binom nr x^r(1-x)^{n-r}\).
Write the unique Bernstein expansion
\[
 1-2^n b_n(x)=\sum_{r=0}^n c_{n,r}\bb_{r,n}(x).
 \tag{1}
\]
Conjecture 3.7 of El Maazouz--Pitman [1, p.~8] proposes positive
coefficients in (1) for every positive \(n\).
\begin{theorem}
Every \(c_{n,r}\) is strictly positive for \(n\geq2\). For \(n=1\),
\((c_{1,0},c_{1,1})=(2,0)\). In particular the expansion has nonnegative
coefficients for every positive \(n\). Moreover, for all \(n\geq1\),
\[
 \max_{0\leq r\leq n}|c_{n,r}-1|
       \leq 4 e^{2\pi}\pi^{-n}.
 \tag{2}
\]
\end{theorem}
\noindent Thus the intended positivity conclusion holds in every degree,
and is strict after the elementary \(n=1\) boundary case.
If \emph{positive} is read as strictly positive even when \(n=1\),
that literal statement has the displayed exception. Restricting the
variable to \([0,1)\) does not change the unique polynomial coefficients.

\section*{A uniform coefficient bound}
The generating function \(te^{xt}/(e^t-1)\) gives
\[
 2^n b_n(x)=\sum_{j=0}^n a_{n,j}x^j,\qquad
 a_{n,j}=\frac{2^n b_{n-j}(0)}{j!}.
\]
Put \(A_\ell=(2\pi)^\ell b_\ell(0)\). Then \(A_0=1\),
\(A_1=-\pi\), \(A_\ell=0\) for odd \(\ell\geq3\), and
\(|A_\ell|=2\zeta(\ell)\leq2\zeta(2)<4\) for positive even \(\ell\).
Consequently
\[
 |a_{n,j}|=\pi^{-n}|A_{n-j}|\frac{(2\pi)^j}{j!}
       \leq4\pi^{-n}\frac{(2\pi)^j}{j!}.
 \tag{3}
\]
\newpage
\section*{Proof for all large degrees and exact small cases}
The monomial-to-Bernstein identity is
\[
 x^j=\sum_{r=0}^n\frac{\binom rj}{\binom nj}\bb_{r,n}(x),
 \qquad 0\leq j\leq n,
 \tag{4}
\]
where \(\binom rj=0\) for \(r<j\).
Indeed, for \(R\sim\operatorname{Binomial}(n,x)\), counting \(j\)-element
sets of successful trials gives
\(\mathbb E\binom Rj=\binom nj x^j\).
The coefficient ratio in (4) belongs to \([0,1]\).
Combining (3)--(4) with \(\sum_r\bb_{r,n}=1\) yields
\[
 |c_{n,r}-1|
 \leq\sum_{j=0}^n|a_{n,j}|
 \leq4\pi^{-n}\sum_{j=0}^\infty\frac{(2\pi)^j}{j!}
 =4e^{2\pi}\pi^{-n}.
\]
This proves (2). For \(n\geq9\), the elementary inequalities
\(3<\pi<7/2\) and \(e<3\) give
\[
 4e^{2\pi}\pi^{-n}<4\cdot3^{7-n}\leq\frac49<1,
\]
so every coefficient in (1) is strictly positive.

For completeness, every remaining coefficient is explicitly certified
below. Formula (4) gives
\[
 c_{n,r}
 =1-\frac{2^n}{n!}\sum_{j=0}^{r}\binom rj B_{n-j}.
 \tag{5}
\]
The table lists a positive denominator \(D_n\) and the entire numerator
row \((D_n c_{n,0},\ldots,D_n c_{n,n})\).
\begin{center}\small
\begin{tabular}{r|r|l}
\(n\)&\(D_n\)&numerator row\\ \hline
1&1&\(2,\ 0\)\\
2&3&\(2,\ 5,\ 2\)\\
3&9&\(9,\ 7,\ 11,\ 9\)\\
4&45&\(46,\ 46,\ 41,\ 46,\ 46\)\\
5&225&\(225,\ 227,\ 229,\ 221,\ 223,\ 225\)\\
6&4725&\(4715,\ 4715,\ 4729,\ 4757,\ 4729,\ 4715,\ 4715\)\\
7&33075&\(33075,\ 33055,\ 33035,\ 33043,\ 33107,\ 33115,\ 33095,\ 33075\)\\
8&33075&\(33082,\ 33082,\ 33077,\ 33067,\ 33059,\ 33067,\ 33077,\ 33082,\ 33082\)
\end{tabular}
\end{center}
These are obtained from (5) using
\[
 (B_0,\ldots,B_8)=
 \left(1,-\tfrac12,\tfrac16,0,-\tfrac1{30},0,
       \tfrac1{42},0,-\tfrac1{30}\right).
\]
All entries for \(2\leq n\leq8\) are positive. For \(n=1\),
\(1-2b_1(x)=2(1-x)=2\bb_{0,1}(x)\), giving the stated boundary
case. This completes the proof in every degree. \(\square\)

Finally, \(\int_0^1\bb_{r,n}(x)\,dx=1/(n+1)\) and
\(\int_0^1b_n(x)\,dx=0\), so
\(\sum_r c_{n,r}=n+1\).
Thus \(c_{n,r}/(n+1)\) are the mixture probabilities when using
normalized beta densities instead of the partition-of-unity basis.
\newpage
\section*{Verification and attribution}
The proof is analytic for every \(n\geq9\), and the complete finite table
certifies all remaining cases without a numerical approximation.
The standard-library verifier independently reconstructs the power
polynomial from the Bernstein coefficients through degree 30, checks
all finite base cases, checks 5,106 further coefficients through degree
100, and checks normalization and reflection through degree 80.
Finite checks supplement the all-degree argument.

The Bernoulli-number formula and the monomial basis conversion are
classical and are not claimed as new. Formula (5) also follows from the
general expansion in Proposition 3.9 of [1]. The result here is a uniform
majorant and its application to the separate degree-\(n\) positivity
target. The approximation theorem for the discrete clock has been
published separately as MathDB 356320.

\section*{Tasks 2--5: four follow-up directions}
\begin{enumerate}[start=2,leftmargin=*]
\item Determine the largest multiplier \(s_n\) for which
\(1-s_n b_n(x)\) has nonnegative degree-\(n\) Bernstein coefficients,
as asked in source Question 3.8.
\item Identify the minimum-coefficient indices for every \(n\), including
their asymptotic location and parity dependence.
\item Find a direct probabilistic construction of the mixture weights
\(c_{n,r}/(n+1)\) using \(n+1\) ordered uniform points.
\item Establish analogous explicit positivity thresholds for other
Appell polynomial densities with uniformly controlled power coefficients.
\end{enumerate}
These directions receive no additional solution credit.

\section*{Dated source and claim audit}
Checked 21 September 2026. All 25 pages of arXiv:2210.02027v2
(11 January 2024) were read during the immediately preceding problem.
The refreshed PDF is byte-identical. Conjecture 3.7 is on page 8.
The author's 2022 thesis repeats it as Conjecture 4.14 on PDF page 89;
only its title and target page were read.
The identified 28-page Cambridge final remains unread after its public
PDF route returned the abstract page.

Initial and final MathDB checks inspect all eight engagement fields.
Recorded bounded claim-index and web searches found no earlier full
target claim in the inspected material, but are not exhaustive.
The source eligibility therefore remains pending and the campaign
count is unchanged. The \(n=1\) wording correction earns no credit.
Review is by the authoring agent, without independent peer review or
proof-assistant certification.

\begin{thebibliography}{1}
\raggedright
\bibitem{emp}
Y. El Maazouz and J. Pitman,
\emph{The Bernoulli clock: probabilistic and combinatorial interpretations
of the Bernoulli polynomials by circular convolution}.
\href{https://arxiv.org/abs/2210.02027}{arXiv:2210.02027v2}, 2024.
Journal final: \emph{Combinatorics, Probability and Computing}
33(2) (2024), 210--237,
\href{https://doi.org/10.1017/S0963548323000421}{doi:10.1017/S0963548323000421}
(identified, full text unread).
\end{thebibliography}
\end{document}
